% GENERATED FILE. Edit canonical lessons, then run npm run generate:books.
% canonical-source-hash: fnv1a64:cbcbaf30fe7d5ee5
% canonical-lessons: PA-C213-request-return, PA-C213-baitho, PA-C213-baitho-review, PA-C213-na-utthho, PA-C213-aje, PA-C213-payoff

\chapter{Please Don't Get Up}
\label{ch:pa-a2-na-request}
\hlchaptermodality{213}

\begin{chapteropening}
\textbf{By the end of this chapter:} \emph{I can ask someone politely to sit or get up, and ask them not to, with na before the request.}

Make a polite request in -o, turn it into a don't with na, and say not yet with aje.
\end{chapteropening}

\section[\texorpdfstring{\pa{ਮਾਫ਼} \pa{ਕਰੋ} · \pa{ਨਾ} (māf karo · nā)}{ਮਾਫ਼ ਕਰੋ · ਨਾ (māf karo · nā)}]{Return to a familiar request}

\label{lesson:PA-C213-request-return}



\begin{quote}
\emph{Hear:} \emph{māf karo} Say \textbf{\pa{ਮਾਫ਼} \pa{ਕਰੋ}}, \textquotedblleft{}excuse me\textquotedblright{}. Then \emph{work with Aman}: \textbf{\pa{ਅਮਨ} \pa{ਨਾਲ} \pa{ਕੰਮ} \pa{ਕਰੋ}}.
\end{quote}

\begin{grammarlens}[title={Grammar Lens}]
Both requests end in \textbf{\pa{ਕਰੋ}}, \textquotedblleft{}do\textquotedblright{}, said to someone politely. You also know \textbf{\pa{ਨਾ}}, the short \textquotedblleft{}no; don't\textquotedblright{}, and two verbs for the body: \textbf{\pa{ਬੈਠਣਾ}}, \textquotedblleft{}to sit\textquotedblright{}, and \textbf{\pa{ਉੱਠਣਾ}}, \textquotedblleft{}to get up\textquotedblright{}.
\end{grammarlens}

\subsection*{Your turn}
\begin{itemize}
\raggedright
  \item \emph{Hear:} \emph{don't} Say \textbf{\pa{ਨਾ}}.
  \item \emph{Hear:} \emph{to sit} Say \textbf{\pa{ਬੈਠਣਾ}}.
  \item \emph{Hear:} \emph{to get up} Say \textbf{\pa{ਉੱਠਣਾ}}.
\end{itemize}

\begin{tcolorbox}[breakable,colback=teal!4,colframe=teal!35!black,title={Before you move on}]
Say \emph{excuse me}, then \emph{don't}. Keep both ready.
\end{tcolorbox}
\section[\texorpdfstring{\pa{ਬੈਠੋ} · \pa{ਉੱਠੋ} (baiṭho · uṭṭho)}{ਬੈਠੋ · ਉੱਠੋ (baiṭho · uṭṭho)}]{\pa{ਬੈਠੋ}, \pa{ਉੱਠੋ} — please sit, please get up}

\label{lesson:PA-C213-baitho}



\begin{quote}
Say \textbf{\pa{ਮਾਫ਼} \pa{ਕਰੋ}}. Then \emph{to sit} and \emph{to get up}.
\end{quote}

\begin{grammarlens}[title={Grammar Lens: stem + \pa{ੋ}}]
\textbf{\pa{ਕਰੋ}} is the stem \textbf{\pa{ਕਰ}} with the sign \textbf{\pa{ੋ}}. Other stems that end in a consonant work the same way:

\begin{quote}
\textbf{\pa{ਬੈਠ}} + \textbf{\pa{ੋ}} $\to$ \textbf{\pa{ਬੈਠੋ}}
\end{quote}

\begin{quote}
\textbf{\pa{ਇੱਥੇ} \pa{ਬੈਠੋ।}} — \emph{itthe baiṭho.} — Please sit here.
\end{quote}

\begin{quote}
\textbf{\pa{ਸਵੇਰੇ} \pa{ਜਲਦੀ} \pa{ਉੱਠੋ।}} — \emph{savere jaldī uṭṭho.} — Get up early in the morning.
\end{quote}

\begin{quote}
\textbf{\pa{ਕਿਤਾਬ} \pa{ਪੜ੍ਹੋ।}} — \emph{kitāb paṛho.} — Read the book.
\end{quote}

It is the polite form, the one you say to \textbf{\pa{ਤੁਸੀਂ}}.
\end{grammarlens}

\subsection*{Your turn}
\begin{itemize}
\raggedright
  \item \emph{Hear:} \emph{please sit there} Say \textbf{\pa{ਉੱਥੇ} \pa{ਬੈਠੋ।}}
  \item \emph{Hear:} \emph{please get up} Say \textbf{\pa{ਉੱਠੋ।}}
  \item \emph{Read it:} \textbf{\pa{ਇੱਥੇ} \pa{ਬੈਠੋ।}} Point to \textbf{\pa{ੋ}} at the end of the request.
  \item \emph{Write it:} copy \textbf{\pa{ਬੈਠੋ}} once
\end{itemize}

\begin{tcolorbox}[breakable,colback=teal!4,colframe=teal!35!black,title={Before you move on}]
\emph{Please sit}? (\textbf{\pa{ਬੈਠੋ}}.) \emph{Please get up}? (\textbf{\pa{ਉੱਠੋ}}.)
\end{tcolorbox}
\section[\texorpdfstring{\pa{ਬੈਠੋ} · \pa{ਉੱਠੋ} (baiṭho · uṭṭho)}{ਬੈਠੋ · ਉੱਠੋ (baiṭho · uṭṭho)}]{Ask politely again}

\label{lesson:PA-C213-baitho-review}



\begin{quote}
\emph{Hear:} \emph{baiṭho; uṭṭho} Say \textbf{\pa{ਬੈਠੋ} · \pa{ਉੱਠੋ}}.
\end{quote}

\subsection*{Your turn}
\begin{itemize}
\raggedright
  \item \emph{Hear:} \emph{please sit here} Say \textbf{\pa{ਇੱਥੇ} \pa{ਬੈਠੋ।}}
  \item \emph{Hear:} \emph{get up early in the morning} Say \textbf{\pa{ਸਵੇਰੇ} \pa{ਜਲਦੀ} \pa{ਉੱਠੋ।}}
  \item \emph{Read it:} \textbf{\pa{ਉੱਠੋ।}} Say \emph{uṭṭho}.
  \item \emph{Write it:} cover the model and write \textbf{\pa{ਉੱਠੋ}} once
\end{itemize}

\begin{tcolorbox}[breakable,colback=teal!4,colframe=teal!35!black,title={Before you move on}]
Ask a guest, politely, to sit.
\end{tcolorbox}
\section[\texorpdfstring{\pa{ਨਾ} \pa{ਉੱਠੋ} (nā uṭṭho)}{ਨਾ ਉੱਠੋ (nā uṭṭho)}]{\pa{ਨਾ} \pa{ਉੱਠੋ} — please don't get up}

\label{lesson:PA-C213-na-utthho}



\begin{quote}
Say \emph{please get up}. Then \emph{don't}, \textbf{\pa{ਨਾ}}. Then \emph{I don't drink tea}.
\end{quote}

\begin{grammarlens}[title={Grammar Lens: \pa{ਨਾ} for a request, \pa{ਨਹੀਂ} for a statement}]
A request takes \textbf{\pa{ਨਾ}}, directly before it:

\begin{quote}
\textbf{\pa{ਨਾ} \pa{ਉੱਠੋ।}} — \emph{nā uṭṭho.} — Please don't get up.
\end{quote}

\begin{quote}
\textbf{\pa{ਉੱਥੇ} \pa{ਨਾ} \pa{ਬੈਠੋ।}} — \emph{utthe nā baiṭho.} — Please don't sit there.
\end{quote}

\begin{quote}
\textbf{\pa{ਹੁਣ} \pa{ਕੰਮ} \pa{ਨਾ} \pa{ਕਰੋ।}} — \emph{huṇ kamm nā karo.} — Don't work now.
\end{quote}

A statement keeps \textbf{\pa{ਨਹੀਂ}}: \textbf{\pa{ਮੈਂ} \pa{ਚਾਹ} \pa{ਨਹੀਂ} \pa{ਪੀਂਦਾ} \pa{ਹਾਂ}}. So the short \textbf{\pa{ਨਾ}} asks, and the long \textbf{\pa{ਨਹੀਂ}} tells.
\end{grammarlens}

\subsection*{Your turn}
\begin{itemize}
\raggedright
  \item \emph{Hear:} \emph{please don't sit there} Say \textbf{\pa{ਉੱਥੇ} \pa{ਨਾ} \pa{ਬੈਠੋ।}}
  \item \emph{Hear:} \emph{don't work now} Say \textbf{\pa{ਹੁਣ} \pa{ਕੰਮ} \pa{ਨਾ} \pa{ਕਰੋ।}}
  \item \emph{Read it:} \textbf{\pa{ਨਾ} \pa{ਉੱਠੋ।}} Give its meaning.
  \item \emph{Write it:} copy \textbf{\pa{ਨਾ} \pa{ਉੱਠੋ}} once
\end{itemize}

\begin{tcolorbox}[breakable,colback=teal!4,colframe=teal!35!black,title={Before you move on}]
\emph{Don't get up}? (\textbf{\pa{ਨਾ} \pa{ਉੱਠੋ}}.) \emph{I don't get up}? (\textbf{\pa{ਮੈਂ} \pa{ਨਹੀਂ} \pa{ਉੱਠਦਾ} \pa{ਹਾਂ}}.)
\end{tcolorbox}
\section[\texorpdfstring{\pa{ਅਜੇ} (aje)}{ਅਜੇ (aje)}]{\pa{ਅਜੇ} (aje) — not yet; still}

\label{lesson:PA-C213-aje}



\begin{quote}
Before the new word: say \emph{please don't get up}, then \emph{now}.
\end{quote}

\subsection*{You'll want to know: \pa{ਅਜੇ}}
\textbf{\pa{ਅਜੇ}} — \emph{aje} — \textquotedblleft{}not yet; still\textquotedblright{}.

\begin{quote}
\textbf{\pa{ਮੈਂ} \pa{ਅਜੇ} \pa{ਘਰ} \pa{ਵਿੱਚ} \pa{ਹਾਂ।}} — \emph{maiṁ aje ghar vichch hāṁ.} — I am still at home.
\end{quote}

\begin{quote}
\textbf{\pa{ਅਜੇ} \pa{ਨਾ} \pa{ਉੱਠੋ।}} — \emph{aje nā uṭṭho.} — Don't get up yet.
\end{quote}

\begin{cousinweb}
\textbf{\pa{ਅਜੇ}} (\emph{aje}) is inherited from Sanskrit \emph{adyaiva}, \textquotedblleft{}even today\textquotedblright{}: CDIAL lists Punjabi \emph{aje}, \textquotedblleft{}hitherto\textquotedblright{} (CDIAL 245). \textquotedblleft{}Up to today\textquotedblright{} is where both of its senses come from: \emph{still}, and with a no, \emph{not yet}.
\end{cousinweb}

\begin{grammarlens}[title={What you've built}]
One more word for saying no and asking.
\end{grammarlens}

\subsection*{Your turn}
\begin{itemize}
\raggedright
  \item \emph{Say it:} \emph{aje}
  \item \emph{Hear:} \emph{I am still at home} Say \textbf{\pa{ਮੈਂ} \pa{ਅਜੇ} \pa{ਘਰ} \pa{ਵਿੱਚ} \pa{ਹਾਂ।}}
  \item \emph{Hear:} \emph{don't get up yet} Say \textbf{\pa{ਅਜੇ} \pa{ਨਾ} \pa{ਉੱਠੋ।}}
  \item \emph{Write it:} copy \textbf{\pa{ਅਜੇ}} once from the model
\end{itemize}

\begin{tcolorbox}[breakable,colback=teal!4,colframe=teal!35!black,title={Before you move on}]
What does \textbf{\pa{ਅਜੇ}} mean? (\textquotedblleft{}Not yet; still\textquotedblright{}.) Say it once more.
\end{tcolorbox}
\section[\texorpdfstring{\pa{ਅਜੇ} \pa{ਨਾ} \pa{ਉੱਠੋ।} (aje nā uṭṭho.)}{ਅਜੇ ਨਾ ਉੱਠੋ। (aje nā uṭṭho.)}]{Please, and please don't: a four-skill payoff}

\label{lesson:PA-C213-payoff}



\begin{quote}
Say \emph{please sit}, then \emph{please don't sit}.
\end{quote}

\subsection*{Your turn: four skills}
\begin{itemize}
\raggedright
  \item \emph{Hear:} \emph{please sit here} Say \textbf{\pa{ਇੱਥੇ} \pa{ਬੈਠੋ।}}
  \item \emph{Hear:} \emph{don't get up yet} Say \textbf{\pa{ਅਜੇ} \pa{ਨਾ} \pa{ਉੱਠੋ।}}
  \item \emph{Read it:} \textbf{\pa{ਉੱਥੇ} \pa{ਨਾ} \pa{ਬੈਠੋ।}} Give its meaning: please don't sit there.
  \item \emph{Write it:} cover the line and write \textbf{\pa{ਅਜੇ} \pa{ਨਾ} \pa{ਉੱਠੋ।}}
  \item Check which no you used: \textbf{\pa{ਨਾ}} for a request, \textbf{\pa{ਨਹੀਂ}} for a statement.
\end{itemize}

\begin{tcolorbox}[breakable,colback=teal!4,colframe=teal!35!black,title={Before you move on}]
A guest stands up to leave too early. Ask them, politely, not to get up yet.
\end{tcolorbox}
